(1) \(E_\pm\) 在 \(\sigma_z\) 基的本征值为 \((1\pm\eta)/2,(1\mp\eta)/2\)，对 \(0\le\eta\le1\) 均非负，且 \(E_++E_-=I\)。结果概率为 \(p_\pm=\operatorname{tr}(\rho E_\pm)=(1\pm\eta z)/2\)。

(2) \(K_+=\operatorname{diag}(\sqrt{(1+\eta)/2},\sqrt{(1-\eta)/2})\)，\(K_-\) 将两对角元交换。归一化选择性态 \(\rho_\pm=K_\pm\rho K_\pm/p_\pm\)，其纵向Bloch分量 \(z_\pm=(z\pm\eta)/(1\pm\eta z)\)；横向分量为 \((x_\pm,y_\pm)=\sqrt{1-\eta^2}(x,y)/(1\pm\eta z)\)。当 \(p_\pm=0\) 时该分支永不出现，不应除以零。

(3) 直接对两个分支不加选择地求和，得到 \(\mathcal E(\rho)_{\uparrow\uparrow}=\rho_{\uparrow\uparrow}\)、\(\mathcal E(\rho)_{\downarrow\downarrow}=\rho_{\downarrow\downarrow}\)，而非对角元满足 \(\mathcal E(\rho)_{\uparrow\downarrow}=\sqrt{1-\eta^2}\rho_{\uparrow\downarrow}\)。因此 \(z\mapsto z\)，\(x,y\mapsto\sqrt{1-\eta^2}(x,y)\)，这是一类可控纯退相干通道。

(4) \(\eta=0\) 时 \(E_\pm=I/2\)，结果与输入无关且通道是恒等操作；\(\eta=1\) 时得到标准 \(\sigma_z\) 投影测量，完全擦除横向相干。两个等先验输入本征态下，以结果符号猜自旋，成功率 \(P_{\rm succ}=(1+\eta)/2\)，信息随 \(\eta\) 增大而相干保留因子 \(\sqrt{1-\eta^2}\) 下降。
